\documentclass[11pt,a4paper]{article}
\usepackage[utf8]{inputenc}
\usepackage[T1]{fontenc}
\usepackage[french]{babel}
\usepackage{geometry}
\usepackage{xcolor}
\usepackage{hyperref}

\geometry{margin=1in}
\definecolor{orange}{RGB}{255,102,0}

\begin{document}

\noindent \textbf{Ismail AISSAOUI} \\
Ingénieur d'État en Intelligence Artificielle (ESI-SBA) \\
Email : ismail.aissaoui@example.com \\
LinkedIn : [Votre Lien]

\bigskip

\begin{flushright}
    À l'attention de l'équipe NAVI / Direction CISS \\
    \textbf{Orange Innovation} \\
    2 Avenue Pierre Marzin \\
    22300 Lannion, France \\
    \medskip
    Le 6 mai 2026
\end{flushright}

\bigskip

\noindent \textbf{Objet : Candidature au doctorat : Apprentissage Supervisé de la Disponibilité des Infrastructures de Virtualisation (Réf : 2026-51771)}

\bigskip

Madame, Monsieur,

Diplômé en tant qu’Ingénieur d’État en Intelligence Artificielle de l’ESI-SBA en juillet 2026, je souhaite soumettre ma candidature pour le poste de doctorat au sein de l’équipe NAVI d’Orange Innovation. Ma motivation pour ce projet réside dans la convergence directe entre mes recherches actuelles sur les Jumeaux Numériques et vos enjeux de résilience des infrastructures Edge Computing.

Actuellement en stage de fin d’études chez **Sonatrach**, je développe un **Jumeau Numérique Génératif** pour une flotte de turbines à gaz LM2500. Ma mission consiste à prédire les états défaillants (Prognostics) en hybridant des lois physiques (**PINNs**) et des architectures séquentielles avancées telles que **Mamba-SSM** et **xLSTM**. Cette approche me permet de capturer des dépendances temporelles complexes tout en garantissant la cohérence physique des prédictions—une problématique que je retrouve dans votre objectif de modélisation de la disponibilité des sites de virtualisation.

Ma maîtrise des outils de la science des données et mon expérience opérationnelle constituent des atouts pour lever les verrous scientifiques de cette thèse :
\begin{itemize}
    \item \textbf{Caractérisation des états dégradés} : Mon expertise en Réseaux de Neurones sur Graphes (**GNN**) est idéale pour modéliser les interdépendances structurelles au sein d'un cluster Kubernetes.
    \item \textbf{Prédiction de défaillance} : J'ai conçu des modèles de calcul de la vie résiduelle (RUL) capables de traiter des flux de données haute fréquence, une compétence directement transposable à votre chaîne de mesure **Prometheus**.
    \item \textbf{Culture du Jumeau Numérique} : Lecteur des travaux de Thelen (cité dans votre offre), je possède une vision claire de l'intégration des modèles prédictifs dans des environnements d'exploitation automatisés.
\end{itemize}

Rejoindre Orange Innovation à Lannion, pôle d'excellence des télécommunications, représenterait pour moi une opportunité unique de transposer mes méthodes d'IA industrielle vers le domaine critique de la virtualisation des réseaux.

Je serais honoré de vous présenter mon parcours et mes propositions techniques lors d'un entretien.

Je vous prie d'agréer, Madame, Monsieur, l'expression de mes salutations distinguées.

\bigskip

\noindent Ismail AISSAOUI

\end{document}
